\section{习题与解答}
\begin{exercise}\label{ex:m4tunnel}
质量$m$的粒子以能量$0<E<V_0$从左射向$0<x<a$的矩形势垒，势垒外$V=0$。求透射率，并给出厚势垒极限，解释它与束缚态边界的联系。
\Diagram{fig:m4tunnelq}{有限禁区两侧都是传播区}{\draw[axis](-1.5,0)--(5,0)node[right]{$x$};\draw[curve](-1.2,0)--(0,0)--(0,2.4)--(3,2.4)--(3,0)--(4.7,0);\draw[guide](-1.2,1)--(4.7,1);\node[right]at(4.7,1){$E$};\node[above]at(1.5,2.4){$V_0$};\node[below]at(3,0){$a$};}
\end{exercise}
\begin{solution}
图\ref{fig:m4tunnelq}取$k=\sqrt{2mE}/\hbar$、$\kappa=\sqrt{2m(V_0-E)}/\hbar$。三段波函数分别为$\ee^{ikx}+r\ee^{-ikx}$、$A\ee^{\kappa x}+B\ee^{-\kappa x}$、$t\ee^{ikx}$。右界面两连续条件给
\[
 A\ee^{\kappa a}=\frac{t\ee^{ika}}2(1+ik/\kappa),\quad
 B\ee^{-\kappa a}=\frac{t\ee^{ika}}2(1-ik/\kappa).
\]
代回左界面，先算出
\begin{align*}
 A+B&=t\ee^{ika}[\cosh\kappa a-i(k/\kappa)\sinh\kappa a],\\
 \kappa(A-B)&=t\ee^{ika}[-\kappa\sinh\kappa a+ik\cosh\kappa a].
\end{align*}
左界面条件为$1+r=A+B$、$ik(1-r)=\kappa(A-B)$。第二式除以$ik$，再与第一式相加，消去$r$，得
\[
 t\ee^{ika}=\left[\cosh\kappa a+
 i\frac{\kappa^2-k^2}{2k\kappa}\sinh\kappa a\right]^{-1}.
\]
\Diagram{fig:m4tunnela}{有限禁区需要两指数支，透射侧只保留向右传播波}{\draw[axis](-1.5,0)--(5,0)node[right]{$x$};\draw[guide](0,-.6)--(0,2.3);\draw[guide](3,-.6)--(3,2.3);\draw[curve,domain=0:3,samples=60]plot(\x,{1.8*exp(-.65*\x)});\draw[->](3.4,.6)--(4.7,.6)node[above]{透射};\node[above]at(1.5,2.4){振幅包络示意};}
两侧传播介质相同，$T=|t|^2$。分母的模平方为
\begin{align*}
 T^{-1}&=\cosh^2\kappa a+
 \frac{(\kappa^2-k^2)^2}{4k^2\kappa^2}\sinh^2\kappa a\\
 &=1+\frac{(\kappa^2+k^2)^2}{4k^2\kappa^2}\sinh^2\kappa a.
\end{align*}
使用$\kappa^2+k^2=2mV_0/\hbar^2$和$k^2\kappa^2=4m^2E(V_0-E)/\hbar^4$，得到
\[
 T=\left[1+\frac{V_0^2\sinh^2\kappa a}{4E(V_0-E)}\right]^{-1}.
\]
图\ref{fig:m4tunnela}的厚垒极限$\kappa a\gg1$给$T\simeq16E(V_0-E)\ee^{-2\kappa a}/V_0^2$。有限禁区内不能先删增长支；它与衰减支一起满足两端匹配，正如有限光隙中的受抑全反射。
\end{solution}

\begin{exercise}\label{ex:m4fieldemission}
金属表面外的电子势垒近似为$V(x)=W-eFx$，$0<x<W/(eF)$，其中$F>0$是外加电场大小。以表面内电子能量为零，求半经典穿透概率的主要指数，解释增大电场的作用。
\Diagram{fig:m4emissionq}{外场使表面势垒成为有限宽三角形}{\draw[axis](-.6,0)--(4.9,0)node[right]{$x$};\draw[axis](0,0)--(0,3.2)node[above]{$V$};\draw[curve](0,2.6)--(3.8,0);\node[left]at(0,2.6){$W$};\node[below]at(3.8,0){$W/(eF)$};\draw[->](-.5,.3)--(-.05,.3);}
\end{exercise}
\begin{solution}
图\ref{fig:m4emissionq}禁区局部虚动量大小$|p|=\sqrt{2m(W-eFx)}$。先计算禁区积分，取$q=W-eFx$，$\dd x=-\dd q/(eF)$，空间上下限对应$q=W,0$：
\begin{align*}
 \int_0^{W/(eF)}|p|\dd x
 &=-\frac{\sqrt{2m}}{eF}\int_W^0q^{1/2}\dd q\\
 &=\frac{\sqrt{2m}}{eF}\left[\frac23q^{3/2}\right]_0^W
 =\frac{2\sqrt{2m}W^{3/2}}{3eF}.
\end{align*}
WKB振幅的衰减因子为$\exp[-\int|p|\dd x/\hbar]$；概率需平方，故
\[
 \mathcal T\propto\exp\left[-\frac2\hbar\int_0^{W/(eF)}|p|\dd x\right]
 =\exp\left[-\frac{4\sqrt{2m}W^{3/2}}{3e\hbar F}\right].
\]
\Diagram{fig:m4emissiona}{电场增大使势垒变窄，指数抑制减弱}{\draw[axis](0,0)--(4.8,0)node[right]{$x$};\draw[axis](0,0)--(0,3.2)node[above]{$V$};\draw[curve](0,2.6)--(4,0);\draw[guide](0,2.6)--(2,0);\node[right]at(3,1.2){较弱场};\node[left]at(-.4,1.4){较强场};\draw[guide](-.3,1.4)--(.8,1.56);}
图\ref{fig:m4emissiona}中势垒高度仍为$W$，宽度却随$1/F$减小，所以穿透概率随电场迅速上升。若画$\ln\mathcal T$对$1/F$，领先指数给负斜率直线；实际电流还含电子供给与前因子，不能仅凭此式求出绝对电流。这个过程不要求入射光子能量超过功函数，和经典过垒光电逃逸属于不同模型。
\end{solution}

\begin{exercise}\label{ex:m4angularvariant}
光电角度模型改为只在向外半球产生电子，概率密度按$\cos\theta$分布，即$\dd P=C\cos\theta\dd\Omega$。法向势垒仍为$W$。求归一化常数和逃逸比例，与全空间各向同性模型比较。
\Diagram{fig:m4angularq}{方向权重偏向表面法线}{\draw[curve](-2,0)--(2,0);\draw[axis](0,0)--(0,2.6)node[above]{法线};\draw[->](0,0)--(.7,2.1);\draw[->](0,0)--(1.7,.8);\node[right]at(1.7,.8){$\cos\theta$权重};}
\end{exercise}
\begin{solution}
图\ref{fig:m4angularq}在向外半球归一化：
\[
 1=2\pi C\int_0^{\pi/2}\cos\theta\sin\theta\dd\theta
 =2\pi C\left[\frac{\sin^2\theta}{2}\right]_0^{\pi/2}
 =\pi C,
\]
故$C=1/\pi$。当$h\nu\geq W$时，可逸出角域为$0\leq\theta\leq\theta_c$，其中$\cos\theta_c=\sqrt{W/(h\nu)}$。在同一原函数中代入新的上限，
\[
 P_{\rm esc}=2\int_0^{\theta_c}\cos\theta\sin\theta\dd\theta
 =[\sin^2\theta]_0^{\theta_c}=1-\frac W{h\nu}.
\]
低于阈值时比例为零。
\Diagram{fig:m4angulara}{同一势垒，不同初始角分布产生不同逃逸曲线}{\draw[axis](0,0)--(4.6,0)node[right]{$h\nu/W$};\draw[axis](0,0)--(0,2.7)node[above]{$P_{\rm esc}$};\draw[curve,domain=1:4,samples=60]plot(\x,{2.3*(1-1/\x)});\draw[guide,domain=1:4,samples=60]plot(\x,{1.15*(1-sqrt(1/\x))});}
图\ref{fig:m4angulara}中此比例在高能趋1，而全空间各向同性模型趋$1/2$。势垒决定允许角域，产生机制决定各方向的权重。因此，相同能量阈值可以对应不同光电产额。
\end{solution}

\begin{exercise}\label{ex:m4radiativecool}
热容量$C$近似恒定的黑球表面积$A$，置于温度$T_0$的大黑腔。球温$T=T_0+\Delta T$，$|\Delta T|\ll T_0$。推导小温差弛豫时间，并说明面积和热容量怎样进入。
\Diagram{fig:m4coolq}{黑球与大腔壁双向辐射交换}{\draw[curve](0,0)circle(.8);\draw[curve](-2,-1.7)rectangle(2,1.7);\node at(0,0){$T$};\node[above]at(0,1.7){$T_0$};\draw[->](.9,.3)--(1.7,.3);\draw[->](-1.7,-.3)--(-.9,-.3);}
\end{exercise}
\begin{solution}
图\ref{fig:m4coolq}中，腔内平衡辐射各向同性，球吸收功率为$A\sigma T_0^4$，自身发射功率为$A\sigma T^4$。以整球热容量$C$列能量方程，得
\[
 C\frac{\dd T}{\dd t}=-A\sigma(T^4-T_0^4).
\]
写$T=T_0+\Delta T$，展开得$T^4-T_0^4=4T_0^3\Delta T+6T_0^2\Delta T^2+\cdots$。在$|\Delta T|/T_0\ll1$时保留线性项，
\[
 \frac{\dd\Delta T}{\dd t}=-\frac{\Delta T}{\tau},\qquad
 \tau=\frac{C}{4A\sigma T_0^3}.
\]
\Diagram{fig:m4coola}{小温差以指数形式衰减}{\draw[axis](0,0)--(4.5,0)node[right]{$t/\tau$};\draw[axis](0,0)--(0,2.6)node[above]{$\Delta T/\Delta T(0)$};\draw[curve,domain=0:4,samples=60]plot(\x,{2.2*exp(-\x)});}
当$\Delta T(0)\ne0$时分离变量，
\[
 \int_{\Delta T(0)}^{\Delta T(t)}\frac{\dd q}{q}
 =-\frac1\tau\int_0^t\dd s,
 \quad \ln\frac{|\Delta T(t)|}{|\Delta T(0)|}=-\frac t\tau.
\]
温差的符号不变，故$\Delta T(t)=\Delta T(0)\ee^{-t/\tau}$，得到图\ref{fig:m4coola}的曲线；初温恰等于腔温时，另有恒解$\Delta T=0$。初温低于腔温时，负温差也按同一式趋于零。量纲上$A\sigma T_0^3$为功率每温度，与$C$相除给时间。

对同材料均匀实心球，$C=\rho c_p(4\pi a^3/3)$、$A=4\pi a^2$，所以$\tau=\rho c_pa/(12\sigma T_0^3)$，半径越大弛豫越慢。该结论要求内部导热足够快，使整球可用单一温度描述。若环境温度可忽略，则原方程改为$\dot T=-(A\sigma/C)T^4$。分离变量，
\[
 \int_{T(0)}^{T(t)}T^{-4}\dd T=-\frac{A\sigma}{C}\int_0^t\dd s,
 \quad \left[-\frac1{3T^3}\right]_{T(0)}^{T(t)}=-\frac{A\sigma t}{C},
\]
得到$T^{-3}(t)=T^{-3}(0)+3A\sigma t/C$，为代数冷却。它采用不同的环境条件，不能沿用小温差指数律。
\end{solution}

\begin{exercise}\label{ex:m4spectralpeak}
求黑体“每对数频率间隔”的能量谱峰值方程，并说明它为什么不同于$u_\nu$的峰。
\Diagram{fig:m4logq}{相同对数宽度在高频处包含更大的线性频率区间}{\draw[axis](0,0)--(5.4,0)node[right]{$\nu$};\draw[very thick](.5,0)--(1,0);\draw[very thick](2,0)--(4,0);\node[above]at(.75,.5){$\nu\to2\nu$};\node[above]at(3,.5){$4\nu\to8\nu$};}
\end{exercise}
\begin{solution}
图\ref{fig:m4logq}比较等比例频率区间。因$\dd\nu=\nu\dd\ln\nu$，对数谱为$\nu u_\nu$，以$x=h\nu/(\kb T)$表示正比$f_4(x)=x^4/(\ee^x-1)$。在例\ref{ex:m4wien}已逐步求出的$f_j'$中取$j=4$，
\[
 f_4'(x)=\frac{x^3\ee^x}{(\ee^x-1)^2}
 [4(1-\ee^{-x})-x].
\]
除去边界根$0$，峰值方程为$4(1-\ee^{-x})=x$，非零正根$x\simeq3.92069$。
\Diagram{fig:m4loga}{对数谱多一个频率权重，峰向高频移动}{\draw[axis](0,0)--(5.3,0)node[right]{$h\nu/(\kb T)$};\draw[axis](0,0)--(0,2.8)node[above]{对数谱};\draw[curve,domain=.05:8,samples=80]plot({.6*\x},{.4*\x^4/(exp(\x)-1)});\draw[guide](2.3524,0)--(2.3524,2);}
图\ref{fig:m4loga}中的峰对应每相同倍频带的能量贡献，而$u_\nu$比较每相同赫兹的能量。它也等于每对数波长间隔能量谱的峰所对应的同一物理频率，因为$|\dd\ln\lambda|=|\dd\ln\nu|$；线性波长谱则还有不同雅可比。
\end{solution}

\begin{exercise}\label{ex:m4welltransition}
无限深方阱宽度缓慢从$L$变为$2L$，粒子始终留在同一量子数$n$的定态分支。求能量与典型动量的变化，并用作用量解释“缓慢”的意义。
\Diagram{fig:m4expandq}{缓慢移动硬壁使阱宽加倍}{\draw[curve](0,2)--(0,0)--(2,0)--(2,2);\draw[guide](4,0)--(4,2);\draw[->](2.2,1)--(3.8,1);\node[below]at(2,-.4){$L\to2L$};}
\end{exercise}
\begin{solution}
图\ref{fig:m4expandq}的定态能量$E_n=n^2h^2/(8mL^2)$，所以阱宽加倍，能量减为$1/4$，动量大小$p_n=nh/(2L)$减为$1/2$。定态内平均动量为0，但$\langle p^2\rangle^{1/2}=p_n$，不要把典型动量大小等同于平均动量。
\Diagram{fig:m4expanda}{相轨道长宽反比变化，面积$2pL=nh$保持}{\draw[axis](0,0)--(5,0)node[right]{$x$};\draw[axis](0,-1.8)--(0,2)node[above]{$p$};\draw[curve](0,-1.3)rectangle(2,1.3);\draw[guide](0,-.65)rectangle(4,.65);}
图\ref{fig:m4expanda}说明缓变过程中作用量$J=2pL$近似保持，量子上要求壁运动足够慢，避免显著跃迁到别的$n$。若壁突然移动，原波函数不再是新阱单一定态，需要展开为新本征态之和，本题能量比例就不能直接使用。
\end{solution}

\begin{exercise}\label{ex:m4oscforce}
谐振子势中加入恒定力$F$，势为$m\omega^2x^2/2-Fx$。求能级、平衡位置和能级间隔，说明零点项是否改变。
\Diagram{fig:m4forceq}{恒力使抛物势横移并下移}{\draw[axis](-2.2,0)--(3.4,0)node[right]{$x$};\draw[axis](0,-1)--(0,3)node[above]{$V$};\draw[guide,domain=-1.8:1.8,samples=50]plot(\x,{.7*\x*\x});\draw[curve,domain=-1:2.8,samples=50]plot(\x,{.7*(\x-.9)^2-.567});}
\end{exercise}
\begin{solution}
逐项求导得$V'(x)=m\omega^2x-F$，令其为零，平衡位置为$x_0=F/(m\omega^2)$；再求导得$V''=m\omega^2>0$，故这是稳定平衡。配方得$V=\tfrac12m\omega^2(x-x_0)^2-F^2/(2m\omega^2)$。图\ref{fig:m4forceq}中的势曲率不变，运动频率仍为$\omega$，能级为$E_n=\hbar\omega(n+1/2)-F^2/(2m\omega^2)$。
\Diagram{fig:m4forcea}{所有能级整体平移，间隔$\hbar\omega$保持}{\foreach\y in{0,.7,1.4,2.1}{\draw[curve](-2,\y)--(-.3,\y);\draw[curve](1,{\y-.5})--(2.7,{\y-.5});}\draw[->](0,1.4)--(.7,.9);\node[above]at(1.8,2.3){有恒力};}
图\ref{fig:m4forcea}中零点能相对势阱最低点仍为$\hbar\omega/2$，但相对原来的能量零点发生共同平移。边界仍是两个平滑转点，故作用量偏移$1/2$不变。只看总能量可能出现负值，并不表示振动激发能或频率为负。
\end{solution}

\begin{exercise}\label{ex:m4hydrogenscaling}
氢样离子核电荷由$Z$变为$2Z$，量子数保持$n$。分别讨论非相对论玻尔半径、速度、束缚能和相邻能级频率的缩放，指出何时这种缩放失效。
\Diagram{fig:m4Zq}{同一$n$下增大核电荷，轨道缩小}{\draw[curve](-1.5,0)circle(1.1);\fill(-1.5,0)circle(2pt);\draw[curve](2,0)circle(.55);\fill(2,0)circle(2pt);\node[below]at(-1.5,-1.3){$Z$};\node[below]at(2,-1.3){$2Z$};}
\end{exercise}
\begin{solution}
图\ref{fig:m4Zq}依$r_n\propto1/Z$、$v_n\propto Z$、$E_n\propto-Z^2$，半径减半、速率加倍、束缚能绝对值乘4。固定两量子数的谱线差同乘4，频率也乘4。
\Diagram{fig:m4Za}{速度尺度$Z\alpha/n$决定非相对论近似的适用性}{\draw[axis](0,0)--(4.5,0)node[right]{$Z\alpha/n$};\draw[axis](0,0)--(0,2.7)node[above]{$v/c$};\draw[curve](0,0)--(3.7,2.3);\node[right]at(3.7,2.3){近1时须改模型};}
图\ref{fig:m4Za}中$Z\alpha/n\ll1$是非相对论条件。增大$Z$会增强相对论修正；当参数逼近1，圆轨道半经典模型本身也失去常规适用性，有限核尺寸、自旋和完整量子动力学不可忽略。缩放律是在固定模型条件下的结果，不是任意核电荷都可无限外推。
\end{solution}

\begin{exercise}\label{ex:m4yield}
采用例\ref{ex:m4escape}的角度筛选模型，在固定吸收功率$P$下改变频率。求零电压光电流最大的光子能量，以及最大电流。
\Diagram{fig:m4yieldq}{光子能量增加时每秒光子数减少，逃逸锥却扩大}{\draw[axis](0,0)--(0,2.7)node[above]{法线};\draw[curve](0,0)--(.8,2);\draw[curve](0,0)--(1.8,2);\node[right]at(1.8,2){较高$h\nu$};\node[left]at(-.1,1){$P$固定};}
\end{exercise}
\begin{solution}
令$x=h\nu/W\geq1$，图\ref{fig:m4yieldq}的逃逸概率为$(1-x^{-1/2})/2$，所以$I=eP( x^{-1}-x^{-3/2})/(2W)$。求导得$\dd I/\dd x=(eP/2W)(-x^{-2}+\tfrac32x^{-5/2})$，零点为$\sqrt x=3/2$。
\Diagram{fig:m4yielda}{电流在$x=9/4$达到极大}{\draw[axis](0,0)--(5,0)node[right]{$x=h\nu/W$};\draw[axis](0,0)--(0,2.6)node[above]{相对电流};\draw[curve,domain=1:7,samples=90]plot({.6*\x},{14*(1/\x-1/(\x^1.5))});\draw[guide](1.35,0)--(1.35,2.075);\node[below]at(1.35,0){$9/4$};}
在$x=1$电流为0，$x\to\infty$也趋于0，故图\ref{fig:m4yielda}的驻点是全局极大。代入得$I_{\max}=2eP/(27W)$、$h\nu=9W/4$。若只看逃逸概率而忽略光子数$P/(h\nu)$，就会错误地认为频率越高电流必越大。
\end{solution}

\begin{exercise}\label{ex:m4cavitypressure}
各向同性平衡辐射能量密度为$u$，理想反射壁法线为$z$。由光子动量通量求辐射压强；再求一束垂直单向光被完全吸收、完全反射的压强。
\Diagram{fig:m4pressureq}{斜向光子反射时法向动量改变$2p\cos\theta$}{\draw[curve](-2,0)--(2,0);\draw[->](-1.6,1.8)--(0,0);\draw[->](0,0)--(1.6,1.8);\draw[guide](0,0)--(0,2.4);\node at(.45,1.3){$\theta$};}
\end{exercise}
\begin{solution}
图\ref{fig:m4pressureq}中一个光子的法向动量改变$2E\cos\theta/c$。单位时间单位面积入射光子的能量还带通量投影$c\cos\theta$，因此
\[
 p=\frac{2u}{4\pi}\int_0^{2\pi}\dd\varphi\int_0^{\pi/2}\cos^2\theta\sin\theta\dd\theta
 =u\left[-\frac{\cos^3\theta}{3}\right]_0^{\pi/2}=\frac u3.
\]
\Diagram{fig:m4pressurea}{单向垂直光束的吸收和反射动量通量}{\draw[curve](0,-.5)--(0,1.2);\draw[->](-2,.8)--(-.1,.8)node[midway,above]{$I/c$};\draw[->](-.1,.2)--(-2,.2)node[midway,below]{反射多一份};\node[right]at(.2,.6){壁};}
图\ref{fig:m4pressurea}的单向光束不需要半球积分，完全吸收时$p=I/c$，完全反射时$p=2I/c$。这里$I$是垂直于光束的功率面密度。腔辐射的$u/3$与孔口通量$cu/4$来自不同角度矩：前者有$\cos^2\theta$，后者有$\cos\theta$。
\end{solution}

\begin{exercise}\label{ex:m4equilibrium}
半径$R$的恒星近似黑体温度$T_s$，距离$r\gg R$处有半径$a$、表面吸收和发射均为黑体的小球。忽略其他热源，求球均匀温度。若只在受光半球均匀发热、背面不辐射，结果如何？
\Diagram{fig:m4equilibriumq}{恒星辐射照到小球的截面积$\pi a^2$}{\draw[curve](0,0)circle(.8);\node at(0,0){$T_s$};\draw[curve](5,0)circle(.35);\foreach\y in{-.2,.2}{\draw[->](1,\y)--(4.6,\y);}\node[below]at(2.8,-.5){$r$};\node[right]at(5.4,0){$a$};}
\end{exercise}
\begin{solution}
恒星总功率$4\pi R^2\sigma T_s^4$，在半径$r$球面上的通量$F=\sigma T_s^4R^2/r^2$。图\ref{fig:m4equilibriumq}的小球吸收面积是迎光截面$\pi a^2$，吸收功率$\pi a^2F$。
\Diagram{fig:m4equilibriuma}{吸收用投影面积，发射用实际表面积}{\draw[curve](0,0)circle(1);\draw[guide](0,-1)--(0,1);\draw[->](1,0)--(1.7,0);\draw[->](0,1)--(0,1.7);\draw[->](-1,0)--(-1.7,0);\draw[->](0,-1)--(0,-1.7);\node[right]at(2,0){全表面$4\pi a^2$};}
均匀温度时，图\ref{fig:m4equilibriuma}全表面发射$4\pi a^2\sigma T^4$。平衡给$T=T_s\sqrt{R/(2r)}$。若只有半球辐射，发射面积$2\pi a^2$，则$T=2^{1/4}T_s\sqrt{R/(2r)}$。该变式是规定的温度分布模型，真实表面是否均温取决于导热、自转等过程。
\end{solution}

\begin{exercise}\label{ex:m4reducedmass}
把玻尔模型的固定核改为质量$M$可动的原子核。求相对坐标的半径、能级及相对固定核结果的修正。
\Diagram{fig:m4reducedq}{电子和核绕共同质心转动}{\draw[curve](0,0)circle(1.7);\draw[curve](0,0)circle(.4);\fill(1.7,0)circle(2pt)node[right]{$m$};\fill(-.4,0)circle(3pt)node[above left]{$M$};\fill(0,0)circle(1pt);\draw[guide](0,0)--(.6,-.65)node[right]{质心};}
\end{exercise}
\begin{solution}
图\ref{fig:m4reducedq}中记电子、原子核位置为$\vct r_e,\vct r_N$，定义质心$\vct R=(m\vct r_e+M\vct r_N)/(m+M)$和相对坐标$\vct r=\vct r_e-\vct r_N$。反解为
\[
 \vct r_e=\vct R+\frac M{M+m}\vct r,\quad
 \vct r_N=\vct R-\frac m{M+m}\vct r.
\]
代入动能，交叉项相消，得到
\[
 K=\tfrac12(M+m)\dot{\vct R}^{\,2}
 +\tfrac12\mu\dot{\vct r}^{\,2},\qquad \mu=\frac{mM}{m+M}.
\]
库仑势仅依赖$r$，因此相对运动等价于质量$\mu$在固定库仑中心中运动。
\Diagram{fig:m4reduceda}{核、电子到质心距离按反质量比分配}{\draw[axis](-1.7,0)--(3.9,0);\fill(-1,0)circle(3pt)node[above]{$M$};\fill(0,0)circle(1.5pt)node[above]{$C$};\fill(3,0)circle(2pt)node[above]{$m$};\draw[<->](-1,-.5)--(0,-.5)node[midway,below]{$r_N$};\draw[<->](0,-.5)--(3,-.5)node[midway,below]{$r_e$};}
在质心系对相对角动量取$\mu vr=n\hbar$，向心方程为$\mu v^2/r=Ze^2/(4\pi\varepsilon_0r^2)$。两式相除先得$v=Z\alpha c/n$，再得
\[
 r_n=\frac{4\pi\varepsilon_0\hbar^2n^2}{\mu Ze^2},\qquad
 E_n=-\frac{\mu c^2(Z\alpha)^2}{2n^2}.
\]
图\ref{fig:m4reduceda}中的电子、核到质心距离分别为$Mr_n/(M+m)$、$mr_n/(M+m)$。相对固定核结果，能级绝对值乘$\mu/m=1/(1+m/M)$，核间距乘其倒数。若$M\gg m$，能级相对修正为$-m/M+O(m^2/M^2)$，相应跃迁频率也降低同一比例。核可动导致的约化质量修正改变内部能级；整个原子发光时的反冲则还需要在这些能级之上作总四动量守恒。
\end{solution}

\begin{exercise}\label{ex:m4powerwell}
粒子在一维势$V(x)=A|x|^s$中运动，$A>0,s>0$。用两平滑转点条件给出能级对量子数的幂次关系，并比较$s=2$与$s\to\infty$的极限。
\Diagram{fig:m4powerq}{幂次束缚势的对称转点}{\draw[axis](-2,0)--(2.3,0)node[right]{$x$};\draw[axis](0,0)--(0,3)node[above]{$V$};\draw[curve,domain=-1.7:1.7,samples=70]plot(\x,{.5*abs(\x)^3});\draw[guide](-1.5,1.6875)--(1.5,1.6875);\node[right]at(1.5,1.6875){$E$};}
\end{exercise}
\begin{solution}
图\ref{fig:m4powerq}中转点为$\pm x_t$，$x_t=(E/A)^{1/s}$。利用势的偶对称，作用量写为
\[
 J=4\int_0^{x_t}\sqrt{2m(E-Ax^s)}\dd x.
\]
令$y=x/x_t$，有$\dd x=x_t\dd y$、$Ax^s=Ey^s$，所以
\[
 J=4\sqrt{2m}A^{-1/s}E^{(s+2)/(2s)}C_s,\quad
 C_s=\int_0^1\sqrt{1-y^s}\dd y.
\]
\Diagram{fig:m4powera}{作用量积分换成固定区间的无量纲面积}{\begin{scope}[yscale=.45]\draw[axis](0,0)--(3.9,0)node[right]{$y$};\draw[axis](0,0)--(0,2.5)node[above]{$\sqrt{1-y^s}$};\draw[curve,domain=0:1,samples=70]plot({3*\x},{2*sqrt(1-\x^3)});\node[below]at(3,0){$1$};\end{scope}}
图\ref{fig:m4powera}的无量纲面积与$E$无关。两软转点条件给出
\[
 E_n\simeq\left[\frac{hA^{1/s}}{4\sqrt{2m}C_s}
 (n+1/2)\right]^{2s/(s+2)}.
\]
因此$s<2$时高能级间距减小，$s=2$时等间隔，$s>2$时增大。取$s=2$，$C_2=\pi/4$；再令$A=m\omega^2/2$，系数化为$\hbar\omega$，恢复谐振子结果。

讨论$s\to\infty$时应先固定长度尺度，写$V=V_*|x/L_*|^s$。对有限$E$，转点趋$L_*$、$C_s\to1$，作用量趋$4L_*\sqrt{2mE}$，与宽$2L_*$的方阱相同，能量的高量子数标度为$n^2$。但端点在极限中变为硬壁，量子数偏移须重新匹配，不能把软转点的$1/2$保留为硬壁最低能级的准确偏移。原点在$s\leq1$时不光滑，转点仍平滑；相关半经典式主要用于高能级，低能级须另解边界问题。
\end{solution}

\begin{exercise}\label{ex:m4gravityscale}
定义重力反弹模型的特征长度$\ell=(\hbar^2/(2m^2g))^{1/3}$。把其薛定谔方程化为无量纲形式，说明精确能量由何种边界条件确定，并比较半经典结果。
\Diagram{fig:m4gravityq}{硬壁处波函数为零，高处波函数衰减}{\draw[axis](-.3,0)--(4.7,0)node[right]{$x$};\node[above]at(0,1.6){$V$};\draw[very thick](0,0)--(0,1.4);\draw[curve](0,0)--(4,1.3);\draw[guide](0,.9)--(4.3,.9);\node[right]at(4.3,.9){$E$};}
\end{exercise}
\begin{solution}
图\ref{fig:m4gravityq}对应$-\hbar^2\psi''/(2m)+mgx\psi=E\psi$。令$\xi=x/\ell$、$\epsilon=E/(mg\ell)$，则链式法则给$\dd\psi/\dd x=\ell^{-1}\dd\psi/\dd\xi$、$\dd^2\psi/\dd x^2=\ell^{-2}\dd^2\psi/\dd\xi^2$。代入并除以$mg\ell$，用$\hbar^2/(2m\ell^2)=mg\ell$，化为
\[
 -\frac{\dd^2\psi}{\dd\xi^2}+\xi\psi=\epsilon\psi.
\]
\Diagram{fig:m4gravitya}{无量纲重力反弹基态在硬壁为零，在高处衰减}{\begin{scope}[yscale=.65]\draw[axis](0,0)--(7.5,0)node[right]{$\xi$};\draw[axis](0,-.4)--(0,2.1)node[above]{$\psi$};\draw[curve]plot coordinates {(0.00000,-0.0000000) (0.06422,0.1348816) (0.12844,0.2685052) (0.19266,0.3996896) (0.25688,0.5273458) (0.32110,0.6504842) (0.38532,0.7682206) (0.44954,0.8797792) (0.51376,0.9844945) (0.57798,1.0818115) (0.64220,1.1712844) (0.70642,1.2525734) (0.77064,1.3254415) (0.83486,1.3897488) (0.89908,1.4454471) (0.96330,1.4925734) (1.02752,1.5312426) (1.09174,1.5616405) (1.15596,1.5840157) (1.22018,1.5986721) (1.28440,1.6059610) (1.34862,1.6062736) (1.41284,1.6000332) (1.47706,1.5876887) (1.54128,1.5697070) (1.60550,1.5465674) (1.66972,1.5187550) (1.73394,1.4867556) (1.79817,1.4510510) (1.86239,1.4121141) (1.92661,1.3704052) (1.99083,1.3263691) (2.05505,1.2804315) (2.11927,1.2329971) (2.18349,1.1844479) (2.24771,1.1351413) (2.31193,1.0854092) (2.37615,1.0355578) (2.44037,0.9858666) (2.50459,0.9365891) (2.56881,0.8879525) (2.63303,0.8401586) (2.69725,0.7933847) (2.76147,0.7477840) (2.82569,0.7034872) (2.88991,0.6606034) (2.95413,0.6192213) (3.01835,0.5794110) (3.08257,0.5412249) (3.14679,0.5046992) (3.21101,0.4698558) (3.27523,0.4367028) (3.33945,0.4052368) (3.40367,0.3754436) (3.46789,0.3472996) (3.53211,0.3207732) (3.59633,0.2958260) (3.66055,0.2724135) (3.72477,0.2504864) (3.78899,0.2299915) (3.85321,0.2108727) (3.91743,0.1930715) (3.98165,0.1765277) (4.04587,0.1611805) (4.11009,0.1469685) (4.17431,0.1338306) (4.23853,0.1217063) (4.30275,0.1105360) (4.36697,0.1002615) (4.43119,0.0908260) (4.49541,0.0821748) (4.55963,0.0742549) (4.62385,0.0670155) (4.68807,0.0604080) (4.75229,0.0543862) (4.81651,0.0489061) (4.88073,0.0439261) (4.94495,0.0394069) (5.00917,0.0353115) (5.07339,0.0316054) (5.13761,0.0282559) (5.20183,0.0252329) (5.26606,0.0225080) (5.33028,0.0200552) (5.39450,0.0178499) (5.45872,0.0158698) (5.52294,0.0140941) (5.58716,0.0125036) (5.65138,0.0110808) (5.71560,0.0098096) (5.77982,0.0086750) (5.84404,0.0076638) (5.90826,0.0067634) (5.97248,0.0059627) (6.03670,0.0052514) (6.10092,0.0046203) (6.16514,0.0040610) (6.22936,0.0035658) (6.29358,0.0031280) (6.35780,0.0027412) (6.42202,0.0023999) (6.48624,0.0020990) (6.55046,0.0018341) (6.61468,0.0016011) (6.67890,0.0013964) (6.74312,0.0012167) (6.80734,0.0010592) (6.87156,0.0009211) (6.93578,0.0008004) (7.00000,0.0006948)};\node[below left]at(0,0){$\psi(0)=0$};\node[above]at(4.6,.6){$\operatorname{Ai}(\xi-\epsilon_0)$};\end{scope}}
令$z=\xi-\epsilon$，方程为$\psi''=z\psi$，通解为$A\operatorname{Ai}(z)+B\operatorname{Bi}(z)$。正无穷处$\operatorname{Bi}$增长，归一化要求$B=0$；$A\ne0$再由归一化确定。图\ref{fig:m4gravitya}的硬壁$x=0$对应$z=-\epsilon$，要求$\operatorname{Ai}(-\epsilon)=0$，所以精确能量由其负零点确定。半经典近似给$\epsilon_n\simeq[\tfrac{3\pi}2(n+\tfrac34)]^{2/3}$。无量纲化把质量和重力全部收进尺度$mg\ell$，可用于比较不同粒子而不重复解微分方程。
\end{solution}
